%=============================================================================
% Chapter 1 Exercises - Rewritten to match original order with hints
%=============================================================================
\section*{Exercises}
\addcontentsline{toc}{section}{Exercises}

\begin{enumerate}
    \item \label{exer1-chapter1}Let $x$ be a nilpotent element of a ring $A$. Show that $1 + x$ is a unit of $A$. Deduce that the sum of a nilpotent element and a unit is a unit.

    \item \label{exer2-chapter1}Let $A$ be a ring and let $A[x]$ be the ring of polynomials in an indeterminate $x$, with coefficients in $A$. Let $f = a_0 + a_1 x + \cdots + a_n x^n \in A[x]$. Prove that:
    \begin{enumerate}[i)]
        \item $f$ is a unit in $A[x]$ $\iff$ $a_0$ is a unit in $A$ and $a_1, \ldots, a_n$ are nilpotent.\\
        \textit{Hint:} If $b_0 + b_1 x + \cdots + b_m x^m$ is the inverse of $f$, prove by induction on $r$ that $a_n^{r+1} b_{m-r} = 0$. Hence show that $a_n$ is nilpotent, and then use Exercise \ref{exer1-chapter1}.
        \item $f$ is nilpotent $\iff$ $a_0, a_1, \ldots, a_n$ are nilpotent.
        \item $f$ is a zero-divisor $\iff$ there exists $a \neq 0$ in $A$ such that $af = 0$.\\
        \textit{Hint:} Choose a polynomial $g = b_0 + b_1 x + \cdots + b_m x^m$ of least degree $m$ such that $fg = 0$. Then $a_n b_m = 0$, hence $a_n g = 0$ (because $a_n g$ annihilates $f$ and has degree $< m$). Now show by induction that $a_{n-r} g = 0$ ($0 \leq r \leq n$).
        \item $f$ is said to be \emph{primitive} if $(a_0, a_1, \ldots, a_n) = (1)$. Prove that if $f, g \in A[x]$, then $fg$ is primitive $\iff$ $f$ and $g$ are primitive.
    \end{enumerate}

    \item \label{exer3-chapter1}Generalize the results of Exercise \ref{exer2-chapter1} to a polynomial ring $A[x_1, \ldots, x_r]$ in several indeterminates.

    \item \label{exer4-chapter1}In the ring $A[x]$, the Jacobson radical is equal to the nilradical.

    \item \label{exer5-chapter1}Let $A$ be a ring and let $A[[x]]$ be the ring of formal power series $f = \sum_{n=0}^\infty a_n x^n$ with coefficients in $A$. Show that:
    \begin{enumerate}[i)]
        \item $f$ is a unit in $A[[x]]$ $\iff$ $a_0$ is a unit in $A$.
        \item If $f$ is nilpotent, then $a_n$ is nilpotent for all $n \geq 0$. Is the converse true? (See Chapter \ref{chapter7}, Exercise \ref{exer2-chapter7}.)
        \item $f$ belongs to the Jacobson radical of $A[[x]]$ $\iff$ $a_0$ belongs to the Jacobson radical of $A$.
        \item The contraction of a maximal ideal $\mathfrak{m}$ of $A[[x]]$ is a maximal ideal of $A$, and $\mathfrak{m}$ is generated by $\mathfrak{m}^c$ and $x$.
        \item Every prime ideal of $A$ is the contraction of a prime ideal of $A[[x]]$.
    \end{enumerate}

    \item \label{exer6-chapter1}A ring $A$ is such that every ideal not contained in the nilradical contains a non-zero idempotent (that is, an element $e$ such that $e^2 = e \neq 0$). Prove that the nilradical and Jacobson radical of $A$ are equal.

    \item \label{exer7-chapter1}Let $A$ be a ring in which every element $x$ satisfies $x^n = x$ for some $n > 1$ (depending on $x$). Show that every prime ideal in $A$ is maximal.

    \item \label{exer8-chapter1}Let $A$ be a ring $\neq 0$. Show that the set of prime ideals of $A$ has minimal elements with respect to inclusion.

    \item \label{exer9-chapter1}Let $\mathfrak{a}$ be an ideal $\neq (1)$ in a ring $A$. Show that $\mathfrak{a} = r(\mathfrak{a})$ $\Rightarrow$ $\mathfrak{a}$ is an intersection of prime ideals.

    \item \label{exer10-chapter1}Let $A$ be a ring, $\nilrad$ its nilradical. Show that the following are equivalent:
    \begin{enumerate}[i)]
        \item $A$ has exactly one prime ideal;
        \item every element of $A$ is either a unit or nilpotent;
        \item $A/\nilrad$ is a field.
    \end{enumerate}

    \item \label{exer11-chapter1}A ring $A$ is \emph{Boolean} if $x^2 = x$ for all $x \in A$. In a Boolean ring $A$, show that:
    \begin{enumerate}[i)]
        \item $2x = 0$ for all $x \in A$;
        \item every prime ideal $\mathfrak{p}$ is maximal, and $A/\mathfrak{p}$ is a field with two elements;
        \item every finitely generated ideal in $A$ is principal.
    \end{enumerate}

    \item \label{exer12-chapter1}A local ring contains no idempotent $\neq 0, 1$.

    \emph{Construction of an algebraic closure of a field}(E. Artin). 
    \item \label{exer13-chapter1}Let $K$ be a field and let $\Sigma$ be the set of all irreducible monic polynomials $f$ in one indeterminate with coefficients in $K$. Let $A$ be the polynomial ring over $K$ generated by indeterminates $x_f$, one for each $f \in \Sigma$. Let $\mathfrak{a}$ be the ideal of $A$ generated by the polynomials $f(x_f)$ for all $f \in \Sigma$. Show that $\mathfrak{a} \neq (1)$.
    
    \quad\, Let $\mathfrak{m}$ be a maximal ideal of $A$ containing $\mathfrak{a}$, and let $K_1 = A/\mathfrak{m}$. Then $K_1$ is an extension field of $K$ in which each $f \in \Sigma$ has a root. Repeat the construction with $K_1$ in place of $K$, obtaining a field $K_2$, and so on. Let $L = \bigcup_{n=1}^\infty K_n$. Then $L$ is a field in which each $f \in \Sigma$ splits completely into linear factors. Let $\bar{K}$ be the set of all elements of $L$ which are algebraic over $K$. Then $\bar{K}$ is an algebraic closure of $K$.

    \item \label{exer14-chapter1}In a ring $A$, let $\Sigma$ be the set of all ideals in which every element is a zero-divisor. Show that the set $\Sigma$ has maximal elements and that every maximal element of $\Sigma$ is a prime ideal. Hence the set of zero-divisors in $A$ is a union of prime ideals.

    \emph{The prime spectrum of a ring.}
    \item \label{exer15-chapter1}Let $A$ be a ring and let $X$ be the set of all prime ideals of $A$. For each subset $E$ of $A$, let $V(E)$ denote the set of all prime ideals of $A$ which contain $E$. Prove that:
    \begin{enumerate}[i)]
        \item if $\mathfrak{a}$ is the ideal generated by $E$, then $V(E) = V(\mathfrak{a}) = V(r(\mathfrak{a}))$.
        \item $V(0) = X$, $V(1) = \varnothing$.
        \item if $(E_i)_{i \in I}$ is any family of subsets of $A$, then $V(\bigcup_{i \in I} E_i) = \bigcap_{i \in I} V(E_i)$.
        \item $V(\mathfrak{a} \cap \mathfrak{b}) = V(\mathfrak{a}\mathfrak{b}) = V(\mathfrak{a}) \cup V(\mathfrak{b})$ for any ideals $\mathfrak{a}, \mathfrak{b}$ of $A$.
    \end{enumerate}
    These results show that the sets $V(E)$ satisfy the axioms for closed sets in a topological space. The resulting topology is called the \emph{Zariski topology}\index{Zariski, ring!topology}. The topological space $X$ is called the \emph{prime spectrum}\index{Spectrum, maximal!prime} of $A$, and is written $\Spec(A)$.

    \item \label{exer16-chapter1}Draw pictures of $\Spec(\Z)$, $\Spec(\R)$, $\Spec(\C[x])$, $\Spec(\R[x])$, $\Spec(\Z[x])$.

    \item \label{exer17-chapter1}For each $f \in A$, let $X_f$ denote the complement of $V(f)$ in $X = \Spec(A)$. The sets $X_f$ are open. Show that they form a basis of open sets for the Zariski topology, and that:
    \begin{enumerate}[i)]
        \item $X_f \cap X_g = X_{fg}$;
        \item $X_f = \varnothing$ $\iff$ $f$ is nilpotent;
        \item $X_f = X$ $\iff$ $f$ is a unit;
        \item $X_f = X_g$ $\iff$ $r((f)) = r((g))$;
        \item $X$ is quasi-compact (that is, every open covering of $X$ has a finite sub-covering).\\
        \textit{Hint:} To prove (v), remark that it is enough to consider a covering of $X$ by basic open sets $X_{f_i}$ ($i \in I$). Show that the $f_i$ generate the unit ideal and hence that there is an equation of the form $1 = \sum_{i \in J} g_i f_i$ ($g_i \in A$) where $J$ is some finite subset of $I$. Then the $X_{f_i}$ ($i \in J$) cover $X$.
        \item More generally, each $X_f$ is quasi-compact.
        \item An open subset of $X$ is quasi-compact if and only if it is a finite union of sets $X_f$.
    \end{enumerate}
    The sets $X_f$ are called \emph{basic open sets} of $X = \Spec(A)$.

    \item \label{exer18-chapter1}For psychological reasons it is sometimes convenient to denote a prime ideal of $A$ by a letter such as $x$ or $y$ when thinking of it as a point of $X = \Spec(A)$. When thinking of $x$ as a prime ideal of $A$, we denote it by $\mathfrak{p}_x$ (logically, of course, it is the same thing). Show that:
    \begin{enumerate}[i)]
        \item the set $\{x\}$ is closed (we say that $x$ is a ``closed point'') in $\Spec(A)$ $\iff$ $\mathfrak{p}_x$ is maximal;
        \item $\overline{\{x\}} = V(\mathfrak{p}_x)$;
        \item $y \in \overline{\{x\}}$ $\iff$ $\mathfrak{p}_x \subseteq \mathfrak{p}_y$;
        \item $X$ is a $T_0$-space (this means that if $x, y$ are distinct points of $X$, then either there is a neighborhood of $x$ which does not contain $y$, or else there is a neighborhood of $y$ which does not contain $x$).
    \end{enumerate}

    \item \label{exer19-chapter1}A topological space $X$ is said to be \emph{irreducible} if $X \neq \varnothing$ and if every pair of non-empty open sets in $X$ intersect, or equivalently if every non-empty open set is dense in $X$. Show that $\Spec(A)$ is irreducible if and only if the nilradical of $A$ is a prime ideal.

    \item \label{exer20-chapter1}Let $X$ be a topological space.
    \begin{enumerate}[i)]
        \item If $Y$ is an irreducible subspace of $X$, then the closure $\bar{Y}$ of $Y$ in $X$ is irreducible.
        \item Every irreducible subspace of $X$ is contained in a maximal irreducible subspace.
        \item The maximal irreducible subspaces of $X$ are closed and cover $X$. They are called the \emph{irreducible components} of $X$. What are the irreducible components of a Hausdorff space?
        \item If $A$ is a ring and $X = \Spec(A)$, then the irreducible components of $X$ are the closed sets $V(\mathfrak{p})$, where $\mathfrak{p}$ is a minimal prime ideal of $A$ (Exercise \ref{exer8-chapter1}).
    \end{enumerate}

    \item \label{exer21-chapter1}Let $\varphi \colon A \to B$ be a ring homomorphism. Let $X = \Spec(A)$ and $Y = \Spec(B)$. If $\mathfrak{q} \in Y$, then $\varphi^{-1}(\mathfrak{q})$ is a prime ideal of $A$, i.e., a point of $X$. Hence $\varphi$ induces a mapping $\varphi^* \colon Y \to X$. Show that:
    \begin{enumerate}[i)]
        \item If $f \in A$ then $\varphi^{*-1}(X_f) = Y_{\varphi(f)}$, and hence that $\varphi^*$ is continuous.
        \item If $\mathfrak{a}$ is an ideal of $A$, then $\varphi^{*-1}(V(\mathfrak{a})) = V(\mathfrak{a}^e)$.
        \item If $\mathfrak{b}$ is an ideal of $B$, then $\varphi^*(V(\mathfrak{b})) = V(\mathfrak{b}^c)$.
        \item If $\varphi$ is surjective, then $\varphi^*$ is a homeomorphism of $Y$ onto the closed subset $V(\Ker(\varphi))$ of $X$. (In particular, $\Spec(A)$ and $\Spec(A/\nilrad)$ (where $\nilrad$ is the nilradical of $A$) are naturally homeomorphic.)
        \item If $\varphi$ is injective, then $\varphi^*(Y)$ is dense in $X$. More precisely, $\varphi^*(Y)$ is dense in $X$ $\iff$ $\Ker(\varphi) \subseteq \nilrad$.
        \item Let $\psi \colon B \to C$ be another ring homomorphism. Then $(\psi\varphi)^* = \varphi^* \circ \psi^*$.
        \item Let $A$ be an integral domain with just one non-zero prime ideal $\mathfrak{p}$, and let $K$ be the field of fractions of $A$. Let $B = (A/\mathfrak{p}) \times K$. Define $\varphi \colon A \to B$ by $\varphi(x) = (\bar{x}, x)$, where $\bar{x}$ is the image of $x$ in $A/\mathfrak{p}$. Show that $\varphi^*$ is bijective but not a homeomorphism.
    \end{enumerate}

    \item \label{exer22-chapter1}Let $A = \prod_{i=1}^n A_i$ be the direct product of rings $A_i$. Show that $\Spec(A)$ is the disjoint union of open (and closed) subspaces $X_i$, where $X_i$ is canonically homeomorphic with $\Spec(A_i)$.
    
    \quad\, Conversely, let $A$ be any ring. Show that the following statements are equivalent:
    \begin{enumerate}[i)]
        \item $X = \Spec(A)$ is disconnected.
        \item $A \cong A_1 \times A_2$ where neither of the rings $A_1, A_2$ is the zero ring.
        \item $A$ contains an idempotent $\neq 0, 1$.
    \end{enumerate}
    In particular, the spectrum of a local ring is always connected (Exercise \ref{exer12-chapter1}).

    \item \label{exer23-chapter1}Let $A$ be a Boolean ring\index{Boolean ring}\index{Ring!Boolean} (Exercise \ref{exer11-chapter1}), and let $X = \Spec(A)$.
    \begin{enumerate}[i)]
        \item For each $f \in A$, the set $X_f$ (Exercise \ref{exer17-chapter1}) is both open and closed in $X$.
        \item Let $f_1, \ldots, f_n \in A$. Show that $X_{f_1} \cup \cdots \cup X_{f_n} = X_f$ for some $f \in A$.
        \item The sets $X_f$ are the only subsets of $X$ which are both open and closed.
        
        \textit{Hint:} Let $Y \subseteq X$ be both open and closed. Since $Y$ is open, it is a union of basic open sets $X_f$. Since $Y$ is closed and $X$ is quasi-compact (Exercise \ref{exer17-chapter1}), $Y$ is quasi-compact. Hence $Y$ is a finite union of basic open sets; now use (ii) above.
        \item $X$ is a compact Hausdorff space.
    \end{enumerate}

    \item \label{exer24-chapter1}Let $L$ be a lattice, in which the sup and inf of two elements $a, b$ are denoted by $a \vee b$ and $a \wedge b$ respectively. $L$ is a \emph{Boolean lattice} (or \emph{Boolean algebra}) if:
    \begin{enumerate}[i)]
        \item $L$ has a least element and a greatest element (denoted by $0, 1$ respectively).
        \item Each of $\vee, \wedge$ is distributive over the other.
        \item Each $a \in L$ has a unique ``complement'' $a' \in L$ such that $a \vee a' = 1$ and $a \wedge a' = 0$.
    \end{enumerate}
    (For example, the set of all subsets of a set, ordered by inclusion, is a Boolean lattice.)

    \quad\, Let $L$ be a Boolean lattice. Define addition and multiplication in $L$ by the rules:
    \[
        a + b = (a \wedge b') \vee (a' \wedge b), \qquad ab = a \wedge b.
    \]
    Verify that in this way $L$ becomes a Boolean ring, say $A(L)$.
    
    \quad\, Conversely, starting from a Boolean ring $A$, define an ordering on $A$ as follows: $a \leq b$ means that $a = ab$. Show that, with respect to this ordering, $A$ is a Boolean lattice. [The sup and inf are given by $a \vee b = a + b + ab$ and $a \wedge b = ab$, and the complement by $a' = 1 - a$.] In this way we obtain a one-to-one correspondence between (isomorphism classes of) Boolean rings and (isomorphism classes of) Boolean lattices.

    \item \label{exer25-chapter1}From the last two exercises deduce \emph{Stone's theorem}, that every Boolean lattice is isomorphic to the lattice of open-and-closed subsets of some compact Hausdorff topological space.

    \item \label{exer26-chapter1}Let $A$ be a ring. The subspace of $\Spec(A)$ consisting of the maximal ideals of $A$, with the induced topology, is called the \emph{maximal spectrum}\index{Spectrum, maximal} of $A$ and is denoted by $\Max(A)$. For arbitrary commutative rings it does not have the nice functorial properties of $\Spec(A)$ (see Exercise \ref{exer21-chapter1}), because the inverse image of a maximal ideal under a ring homomorphism need not be maximal.
    
    \quad\, Let $X$ be a compact Hausdorff space and let $C(X)$ denote the ring of all real-valued continuous functions on $X$ (add and multiply functions by adding and multiplying their values). For each $x \in X$, let $\mathfrak{m}_x$ be the set of all $f \in C(X)$ such that $f(x) = 0$. The ideal $\mathfrak{m}_x$ is maximal, because it is the kernel of the (surjective) homomorphism $C(X) \to \R$ which takes $f$ to $f(x)$. If $\mathfrak{X}$ denotes $\Max(C(X))$, we have therefore defined a mapping $\mu \colon X \to \mathfrak{X}$, namely $x \mapsto \mathfrak{m}_x$.

    \quad\, We shall show that $\mu$ is a homeomorphism of $X$ onto $\mathfrak{X}$.
    \begin{enumerate}[i)]
        \item Let $\mathfrak{m}$ be any maximal ideal of $C(X)$, and let $V = V(\mathfrak{m})$ be the set of common zeros of the functions in $\mathfrak{m}$: that is, $V = \{x \in X : f(x) = 0 \text{ for all } f \in \mathfrak{m}\}$. Suppose that $V$ is empty. Then for each $x \in X$ there exists $f_x \in \mathfrak{m}$ such that $f_x(x) \neq 0$. Since $f_x$ is continuous, there is an open neighborhood $U_x$ of $x$ in $X$ on which $f_x$ does not vanish. By compactness a finite number of the neighborhoods, say $U_{x_1}, \ldots, U_{x_n}$, cover $X$. Let $f = f_{x_1}^2 + \cdots + f_{x_n}^2$. Then $f$ does not vanish at any point of $X$, hence is a unit in $C(X)$. But this contradicts $f \in \mathfrak{m}$, hence $V$ is not empty. Let $x$ be a point of $V$. Then $\mathfrak{m} \subseteq \mathfrak{m}_x$, hence $\mathfrak{m} = \mathfrak{m}_x$ because $\mathfrak{m}$ is maximal. Hence $\mu$ is surjective.
        \item By Urysohn's lemma (this is the only non-trivial fact required in the argument) the continuous functions separate the points of $X$. Hence $x \neq y$ $\Rightarrow$ $\mathfrak{m}_x \neq \mathfrak{m}_y$, and therefore $\mu$ is injective.
        \item Let $f \in C(X)$; let $U_f = \{x \in X : f(x) \neq 0\}$ and let $\mathcal{U}_f = \{\mathfrak{m} \in \mathfrak{X} : f \notin \mathfrak{m}\}$. Show that $\mu(U_f) = \mathcal{U}_f$. The open sets $U_f$ (resp. $\mathcal{U}_f$) form a basis of the topology of $X$ (resp. $\mathfrak{X}$) and therefore $\mu$ is a homeomorphism.
        
        Thus $X$ can be reconstructed from the ring of functions $C(X)$.
    \end{enumerate}
    

    \emph{Affine algebraic varieties.}\index{Varieties, affine algebraic}
    \item \label{exer27-chapter1} Let $k$ be an algebraically closed field and let $\{f_\alpha(t_1, \ldots, t_n) = 0\}$ be a set of polynomial equations in $n$ variables with coefficients in $k$. The set $X$ of all points $x = (x_1, \ldots, x_n) \in k^n$ which satisfy these equations is an affine algebraic variety.

    Consider the set of all polynomials $g \in k[t_1, \ldots, t_n]$ with the property that $g(x) = 0$ for all $x \in X$. This set is an ideal $I(X)$ in the polynomial ring, and is called the \emph{ideal of the variety} $X$. The quotient ring $P(X) = k[t_1, \ldots, t_n]/I(X)$ is the ring of polynomial functions on $X$, because two polynomials $g, h$ define the same polynomial function on $X$ if and only if $g - h$ vanishes at every point of $X$, that is, if and only if $g - h \in I(X)$.
    
    \quad\, Let $\xi_i$ be the image of $t_i$ in $P(X)$. The $\xi_i$ ($1 \leq i \leq n$) are the coordinate functions on $X$: if $x \in X$, then $\xi_i(x)$ is the $i$th coordinate of $x$. $P(X)$ is generated as a $k$-algebra by the coordinate functions, and is called the \emph{coordinate ring} (or \emph{affine algebra}) of $X$.

    \quad\, As in Exercise \ref{exer26-chapter1}, for each $x \in X$ let $\mathfrak{m}_x$ be the ideal of all $f \in P(X)$ such that $f(x) = 0$; it is a maximal ideal of $P(X)$. Hence, if $\tilde{X} = \Max(P(X))$, we have defined a mapping $\mu \colon X \to \tilde{X}$, namely $x \mapsto \mathfrak{m}_x$.

    It is easy to show that $\mu$ is injective: if $x \neq y$, we must have $x_i \neq y_i$ for some $i$ ($1 \leq i \leq n$), and hence $\xi_i - x_i$ is in $\mathfrak{m}_x$ but not in $\mathfrak{m}_y$, so that $\mathfrak{m}_x \neq \mathfrak{m}_y$. What is less obvious (but still true) is that $\mu$ is surjective. This is one form of \emph{Hilbert's Nullstellensatz} (see Chapter \ref{chapter7}).

    \item \label{exer28-chapter1}Let $f_1, \ldots, f_m$ be elements of $k[t_1, \ldots, t_n]$. They determine a polynomial mapping $\varphi \colon k^n \to k^m$: if $x \in k^n$, the coordinates of $\varphi(x)$ are $f_1(x), \ldots, f_m(x)$.

    \quad\, Let $X, Y$ be affine algebraic varieties in $k^n, k^m$ respectively. A mapping $\varphi \colon X \to Y$ is said to be \emph{regular} if $\varphi$ is the restriction to $X$ of a polynomial mapping from $k^n$ to $k^m$.

    \quad\, If $\eta$ is a polynomial function on $Y$, then $\eta \circ \varphi$ is a polynomial function on $X$. Hence $\varphi$ induces a $k$-algebra homomorphism $P(Y) \to P(X)$, namely $\eta \mapsto \eta \circ \varphi$. Show that in this way we obtain a one-to-one correspondence between the regular mappings $X \to Y$ and the $k$-algebra homomorphisms $P(Y) \to P(X)$.
\end{enumerate}
